\section{引言}

QCD 最重要的目标之一，是从其基本自由度——夸克与胶子——出发理解强子结构。这必然进入 QCD 的非微扰区，一般而言，难以直接从 QCD 拉氏量获得第一性原理的结果。格点 QCD 是获得此类结果的有力工具，它是一种定义在欧氏时空上的方法，已被用于以惊人的精度计算强子质量、电荷等物理量。然而，它无法用来直接触及本质上属于闵可夫斯基时空的观测量，例如部分子分布函数（PDF）。PDF 定义为光锥关联的前向强子矩阵元，描述强子内部夸克与胶子的动量分布；它们在理解诸如 LHC 这类高能强子对撞机的实验数据方面起着关键作用。由于其含实时依赖性，格点 QCD 只能通过计算各阶矩来间接提取 PDF 的信息，而这受到诸如算符混合等技术复杂性的限制。另一种常用的确定 PDF 的方法，是假定一个合适的参数化形式并拟合大量各式各样的实验数据。多数 PDF 合作组即以此方式得到其 PDF 集合~\cite{Alekhin:2012ig,Gao:2013xoa,Radescu:2010zz,CooperSarkar:2011aa,Martin:2009iq,Ball:2012cx}。该方法的一个主要缺点是存在参数化不确定性，而且不同合作组对同一 PDF 往往给出不同的结果。


近期的发展~\cite{Ji:2013fga,Ji:2013dva,Xiong:2013bka,Lin:2014zya,Hatta:2013gta,Ma:2014jla,Ji:2014hxa,Ji:2014gla,Ji:2014lra,
Alexandrou:2015rja,Ji:2015jwa,Ji:2015qla,Xiong:2015nua,Li:2016amo,Chen:2016utp} 表明，诸如 PDF 这样的光锥观测量可以由一个类空关联子的强子矩阵元在大动量极限下直接提取；该矩阵元被称为准观测量（quasi observable），所用的框架是大动量有效理论~\cite{Ji:2014lra}（关于提取光锥量的其他途径，参见例如文献~\cite{Davoudi:2012ya,Detmold:2005gg,Liu:1993cv,Liu:1998um,Liu:1999ak,Liu:2016djw,Braun:2007wv}）。准 PDF 不具有实时依赖性，因而可以在格点上模拟。通过直接计算已证明：味道非单态准 PDF 与其对应的光锥 PDF 在单圈层次具有相同的红外（IR）行为；文献~\cite{Xiong:2013bka} 进一步论证了这一点在所有圈次下均成立，并据此给出了因子化公式。

文献~\cite{Xiong:2013bka} 中的因子化是针对裸准夸克分布给出的，即进入准分布的所有场与耦合都是裸量。然而，与光锥分布一样，准分布也包含紫外（UV）发散，因而需要重正化。文献~\cite{Ji:2015jwa} 探讨了准分布的重正化性质，并证明它在两圈层次可乘法重正化。该文还在准夸克分布的虚修正与重夸克有效理论中重轻夸克矢量流的修正之间建立了等价关系，从而使前者的 UV 发散可以按后者的重正化方式处理。文献~\cite{Ji:2015jwa} 采用的是维数正规化，截断或格点正规化中出现的线性发散被忽略了。在现实的格点计算中，必须知道如何处理这类幂次发散。这是本文的目标之一。我们将证明：准夸克分布中的幂次发散可以通过一个质量反项在所有圈次上消除，这与开放威尔逊线的重正化完全相同。经过这样的质量重正化之后，准夸克分布得到改进，至多只含对数发散。

本文其余部分安排如下：第 2 节讨论准夸克分布中威尔逊线自能导致的幂次发散的重正化；第 3 节给出准夸克分布的格点微扰论匹配；最后在第 4 节给出结论。
